%------------------------------------------------------------------------------%
% Luis A. D'Afonseca
%
%------------------------------------------------------------------------------%
\documentclass[a4paper,12pt,fleqn]{article}

\usepackage{style_avaliacao}

\ShowAnswers

%------------------------------------------------------------------------------%
\begin{document}

\makeHeader{Integrais e Séries}{Prova 4 -- Turma 4}{12/12/2023}

% 01
%------------------------------------------------------------------------------%
\exercise{30}
Calcule a Série de Taylor da função
\(
  f(x) = \int \sinh\left(x^2\right) dx
\)

\textit{Lembrete}: \(\sinh(x)=\frac{e^x-e^{-x}}{2}\)

\begin{answer}
  Sabemos que a Série de Taylor da função exponencial é
  \[
    e^x = \sum_{n=0}^\infty \frac{x^n}{n!}
  \]
  Primeiro calculamos a Série de Taylor das funções
  \[
    e^{x^2}
    = \sum_{n=0}^\infty \frac{\left(x^2\right)^n}{n!}
    = \sum_{n=0}^\infty \frac{x^{2n}}{n!}
  \]
  \[
    e^{-x^2}
    = \sum_{n=0}^\infty \frac{\left(-x^2\right)^n}{n!}
    = \sum_{n=0}^\infty \frac{(-1)^nx^{2n}}{n!}
  \]
  agora da função
  \begin{align*}
    \sinh\left(x^2\right)
    & = \frac{e^{x^2}-e^{-x^2}}{2} \\[2mm]
    & = \frac{1}{2}\left(
          \sum_{n=0}^\infty \frac{       x^{2n}}{n!} -
          \sum_{n=0}^\infty \frac{(-1)^n x^{2n}}{n!}
        \right) \\[2mm]
    & = \frac{1}{2} \sum_{n=0}^\infty \frac{1}{n!}
                    \big[x^{2n}-(-1)^n x^{2n}\big] \\[2mm]
    & = \frac{1}{2} \sum_{n=0}^\infty \frac{1-(-1)^n}{n!} x^{2n}
  \end{align*}
  Observando que
  \(
    1 - (-1)^n = 1 + (-1)^{n+1}
  \)
  vale 2 para $n$ ímpar e zero para $n$ par,
  podemos remover todos os índices pares e
  fazer a mudança de índice \(n=2k+1\), obtendo
  \[
    \sinh\left(x^2\right)
    = \frac{1}{2} \sum_{k=0}^\infty \frac{2}{(2k+1)!} x^{2(2k+1)}
    = \sum_{k=0}^\infty \frac{x^{4k+2}}{(2k+1)!}
  \]
  Agora integramos
  \begin{align*}
    f(x)
    & = \int \sinh\left(x^2\right) dx \\[2mm]
    & = \int \sum_{k=0}^\infty \frac{x^{4k+2}}{(2k+1)!} \;dx \\[2mm]
    & = \sum_{k=0}^\infty \frac{1}{(2k+1)!} \int x^{4k+2} dx \\[2mm]
    & = C + \sum_{k=0}^\infty \frac{1}{(2k+1)!} \frac{x^{4k+3}}{4k+3} \\[2mm]
    & = C + \sum_{k=0}^\infty \frac{x^{4k+3}}{(4k+3)(2k+1)!}
  \end{align*}
\end{answer}

% 02
%------------------------------------------------------------------------------%
\exercise{40}
Seja \( f(x) = \frac{2}{1+x}\)
\begin{enumerate}
  \item[{a) [12]}] Encontre um padrão para \(f^{(n)}(x)\)
  \item[{b) [14]}] Usando o padrão do item anterior, encontre a
                   Serie de Taylor de $f(x)$, centrada em zero
  \item[{c) [14]}] Obtenha o intervalo de convergência da série de potências
\end{enumerate}

\begin{answer}
  \begin{enumerate}[label=\alph*)]
    \item Calculando as primeiras derivadas
    \begin{align*}
      f^{(0)}(x) & = \frac{2}{1+x} = 2(1+x)^{-1} \\[2mm]
      f^{(1)}(x) & = \frac{\,d\;}{\,dx\;}\left[2(1+x)^{-1}\right]
      = 2(-1)(1+x)^{-1-1} (1+x)'
      = -2 (1+x)^{-2} \\[2mm]
      f^{(2)}(x) & = \frac{d^2}{dx^2}\left[2(1+x)^{-1}\right]
      = \hpm 2 \cdot 2                 (1+x)^{-3} \\[2mm]
      f^{(3)}(x) & = \frac{d^3}{dx^3}\left[2(1+x)^{-1}\right]
      = -2 \cdot 2 \cdot 3         (1+x)^{-4} \\[2mm]
      f^{(4)}(x) & = \frac{d^4}{dx^4}\left[2(1+x)^{-1}\right]
      = \hpm 2 \cdot 2 \cdot 3 \cdot 4 (1+x)^{-5} \\[2mm]
      f^{(5)}(x) & = \frac{d^5}{dx^5}\left[2(1+x)^{-1}\right]
      = - 2 \cdot 2 \cdot 3 \cdot 4 \cdot 5 (1+x)^{-6}
    \end{align*}
    Por inspeção, vemos que
    \[
      f^{(n)}(x)
      = \frac{d^n}{dx^n}\left[2(1+x)^{-1}\right]
      = 2(-1)^n n! (1+x)^{-n-1}
      = \frac{2 (-1)^n n!}{(1+x)^{n+1}}
    \]
    \item Os coeficientes da série de Taylor são
    \[
      c_n
      = \frac{f^{(n)}(0)}{n!}
      = \frac{2(-1)^n n!}{(1-0)^{n+1}n!}
      = 2(-1)^n
    \]
    Montando a série temos
    \[
      T(x)
      = \sum_{n=0}^\infty c_n x^n
      = \sum_{n=0}^\infty 2(-1)^n x^n
    \]
    \item Reescrevendo a série temos
    \[
      T(x)
      = \sum_{n=0}^\infty 2 (-1)^n x^n
      = \sum_{n=0}^\infty 2 (-x)^n
    \]
    que é uma série geométrica com \(\alpha = 2\) e \(r = -x\),
    portanto, convergente apenas quando \(\abs{r}<1\) ou seja \(x\in(-1, 1)\).
  \end{enumerate}
\end{answer}

% 03
%------------------------------------------------------------------------------%
\exercise{30}
Prove que a série de Maclaurin da função cosseno
converge para $\cos x$, em todo $x\in\R$

\begin{answer}
  Pelo \textbf{Teorema de Taylor}, sabemos que o erro cometido ao aproximarmos
  $\cos(x)$ pelo seu polinômio de Taylor centrado em zero de grau $n$ é
  \[
    R_n(x) = \frac{\cos^{(n+1)}(\xi)}{(n+1)!} \, x^{n+1}
  \]
  para algum $\xi$ entre zero e $x$.
  Como todas as derivadas do cosseno são cossenos ou senos alternando os sinais,
  podemos escrever
  \[
    \abs{\cos^{(k)}(t)} \leq 1
  \]
  para todo \(k\in\N\) e \(t\in\R\).
  Portanto, para todo \(x\in\R\)
  \[
    \abs{R_n(x)}
    = \frac{\abs{\cos^{(n+1)}(\xi)}}{(n+1)!}\abs{x^{n+1}}
    \leq \frac{\abs{x^{n+1}}}{(n+1)!}
  \]
  Calculando o limite
  \[
    \lim_{n\to\infty} \frac{\abs{x^{n+1}}}{(n+1)!}
    = \abs{x^{n+1}} \lim_{n\to\infty} \frac{1}{(n+1)!}
    = \abs{x^{n+1}} \cdot 0
    = 0
  \]
  podemos usar o \textbf{Teorema do Confronto} para mostrar que \(R_n(x)\)
  vai para zero, quando $n$ vai para o infinito.
  Portanto, a série converge para a função.
\end{answer}

%------------------------------------------------------------------------------%
\end{document}
%------------------------------------------------------------------------------%
